% Esquema vectorial basado en Ray RLlib, Key Concepts (consulta: 2026-10-26).
% Cada flecha representa datos o parámetros; no una asignación de hardware.
\begingroup
\setstretch{1}
\begin{tikzpicture}[met/base,
 rllib/group/.style={draw=black!50,rounded corners=3pt,line width=.7pt},
 rllib/node/.style={met/box,text width=4.8cm,minimum height=.85cm},
 rllib/label/.style={met/note,fill=white,inner sep=2pt}]
\draw[rllib/group,fill=colEntorno!8] (0,0) rectangle (6,-5.2);
\draw[rllib/group,fill=colInterfaz!10] (8,0) rectangle (14,-5.2);
\node[met/head] at (3,-.40) {Muestreo};
\node[met/head] at (11,-.40) {Aprendizaje};
\node[met/note] at (3,-.88) {\texttt{EnvRunnerGroup}};
\node[met/note] at (11,-.88) {\texttt{LearnerGroup}};
\node[rllib/node] (env) at (3,-2.0) {Entorno de simulación};
\node[rllib/node] (inference) at (3,-4.1)
 {\texttt{RLModule}\\Elección de acciones};
\node[rllib/node,fill=colInterfaz!20] (update) at (11,-2.0)
 {\textit{Learner}\\Pérdidas y gradientes\\Optimizador};
\node[rllib/node,fill=colInterfaz!20] (model) at (11,-4.1)
 {\texttt{RLModule}\\Pesos del modelo};
\draw[met/arrow] ([xshift=-2cm]env.south) --
 node[rllib/label,right,text width=2.3cm] {Observaciones\\y recompensas}
 ([xshift=-2cm]inference.north);
\draw[met/arrow] ([xshift=1.1cm]inference.north) --
 node[rllib/label,right] {Acciones} ([xshift=1.1cm]env.south);
\draw[met/arrow] (update.south) --
 node[rllib/label,right] {Actualización} (model.north);
\draw[met/arrow] (env.east) --
 node[rllib/label,above,text width=2.0cm] {Datos de\\entrenamiento}
 (update.west);
\draw[met/arrow] (model.south) -- (11,-5.75) --
 node[rllib/label,below] {Sincronización de pesos}
 (3,-5.75) -- (inference.south);
\end{tikzpicture}
\endgroup
